%% Separately authored historical derivative following exact source lines 10675--11353.
%% Necessity: exact restoration replaces the historical constructibility and
%% forcing rewrite, but does not contain the distinct condensation theorem,
%% apparatus A-131/A-132, or ccc/name-counting consequences retained here.
%% Definition: this file is based on historical body slices 87--128, 296--323,
%% and 324--372; consumers are redirected to exact source, caption IDs are
%% cleaned, and the later forcing scope is explicitly renewed.
%% Standing: this is separately authored reader material, not inherited source
%% or a promotion of the exact forcing theorem.
%% Nonredundancy: the exact source supplies all constructibility and forcing core
%% statements/proofs; these slices supply only additional cited/appartus results.
%% Use: parts/p06.tex routes this file once after all sixteen exact wrappers, and
%% its three proved results have reciprocal entries in the bounded annex derivative.

\section{Further constructibility and forcing apparatus}

% BEGIN HISTORICAL SLICE 87-128
\begin{hbtheorem}{thm:p06-condensation}{cited=Jech2003}{Condensation for \(L\)}
Let \(\alpha>\omega\) be a limit ordinal and let
  \(X\preccurlyeq(L_\alpha,\in)\) be a set-sized elementary substructure.  The
Mostowski collapse of \((X,\in)\) is \((L_\beta,\in)\) for some
\(\beta\leq\alpha\).  Elementarity is load-bearing: arbitrary subsets of
\(L_\alpha\) do not satisfy this conclusion.
\end{hbtheorem}

Jensen's fine-structure analysis is a further-reading pointer rather than a
local entry; see \cite{Devlin1984}.  The finite-family L\'evy reflection statement is
Theorem~\ref{thm:p06-levy-reflection}; its formula-family and parameter
hypotheses are stated there rather than weakened here.

\begin{figure}[htbp]
\centering
\begin{tikzpicture}[scale=.85,transform shape,>=Stealth,
  level/.style={draw,rounded corners,minimum width=3.6cm,
    minimum height=.62cm,fill=blue!5}]
  \node[level] (l0) at (0,0) {$L_0=\varnothing$};
  \node[level] (l1) at (0,1.05)
    {$L_{\alpha+1}=\operatorname{Def}(L_\alpha)$};
  \node[level] (ll) at (0,2.10)
    {$L_\lambda=\bigcup_{\beta<\lambda}L_\beta$};
  \node[level,minimum width=5.3cm] (linf) at (0,3.15)
    {$L=\bigcup_{\alpha\in\operatorname{Ord}}L_\alpha\subseteq V$};
  \draw[->] (l0) -- (l1);
  \draw[->] (l1) -- (ll);
  \draw[->] (ll) -- (linf);
  \node[align=left,anchor=west] at (3.1,1.58)
    {successor stages retain only\\subsets definable over the prior level};
  \draw[->,dashed] (2.95,1.58) -- (1.85,1.15);
  \node[align=center] at (0,4.0)
    {Condensation identifies suitable elementary substructures\\
     with earlier levels; \(V=L\) asserts \(V=L\).};
\end{tikzpicture}
\caption{The constructible hierarchy inside the cumulative universe:
each \(L_\alpha\subseteq V_\alpha\), and the union defining \(L\) continues
through every ordinal, so \(L\) is a proper class.}
\label{fig:p06-constructible-hierarchy}
\end{figure}

% END HISTORICAL SLICE 87-128

% BEGIN HISTORICAL SLICE 296-323
\begin{figure}[htbp]
\centering
\begin{tikzpicture}[scale=.74,transform shape,>=Stealth,node distance=2.4cm,
  box/.style={draw,rounded corners,minimum width=3.0cm,
    minimum height=1.05cm,align=center}]
  \node[box,fill=blue!5] (m) {\(M\)\\names \(\nu\), poset \(\mathbb P\)};
  \node[box,fill=green!7,right=of m] (g)
    {external \(M\)-generic\\filter \(G\subseteq\mathbb P\)\\
     \(G\notin M\) in the Cohen case};
  \node[box,fill=orange!8,right=of g] (mg)
    {\(M[G]\)\\values \(\nu^G\)};
  \node[box,fill=yellow!8,below=1.0cm of g] (ev)
    {name evaluation\\\((\nu,G)\mapsto\nu^G\)};
  \draw[->] (m) -- node[above]{dense sets} (g);
  \draw[->] (m.south east) -- node[below left]{names} (ev.west);
  \draw[->] (g.south) -- node[right]{parameter} (ev.north);
  \draw[->] (ev.east) -- node[below right]{values} (mg.south west);
  \draw[->,bend left=25] (m) to
    node[above]{\(x\mapsto\check x\mapsto x\)} (mg);
  \node[below=1.0cm of ev,align=center]
    {The pedagogical construction assumes external countability.\\
     Cited forcing theorem: \(M[G]\entails\mathsf{ZFC}\) and
     \(\operatorname{Ord}^{M[G]}=\operatorname{Ord}^M\).};
\end{tikzpicture}
\caption{Names, a generic filter, and the forcing extension.}
\label{fig:p06-forcing-extension}
\end{figure}

% END HISTORICAL SLICE 296-323

% BEGIN HISTORICAL SLICE 324-372
\section{Chain condition and many Cohen reals}
\label{sec:p06-ccc-cohen}

\noindent\textit{Reader note.}  This separately authored section supplements
the exact forcing and relative-consistency sequence.  It does not replace or
upgrade the exact semantic-forcing definition or the cited forcing theorem.

\begin{hbconvention}{def:p06-ccc-forcing-standing}{Standing context for ccc and name counting}
Throughout this section, \(M\) is the externally countable transitive set
model of \(\mathsf{ZFC}\) fixed in
Definition~\ref{def:p06-forcing-order-orientation},
\(\mathbb P\in M\) is a forcing poset with the same stronger-is-smaller
orientation, and \(G\subseteq\mathbb P\) is an external
\(M\)-generic poset filter.  Cardinalities and exponentiation carrying the
superscript \(M\) are computed in \(M\).

For each fixed standard formula and tuple of names, the bare notation
\(p\forces\varphi(\bar\nu)\) abbreviates the recursively defined assertion
imported by the cited forcing theorem.  The exact indexed notation
\(p\forces^{\mathrm{sem}}_{M,\mathbb P}\varphi(\bar\nu)\) remains the external
semantic relation.  The forcing theorem states their equivalence; this
convention imports that result and does not claim a second local recursion.
The distinction is used by the truth-lemma and nice-name arguments below.
\end{hbconvention}

\begin{hbtheorem}{thm:p06-ccc-preservation}{proved=annex:p6.force.9}{Cardinal and cofinality preservation by ccc forcing}
Assume \((M,\in)\entails\mathsf{ZFC}\), \(M\) regards \(\mathbb P\) as ccc,
and \(G\) is \(M\)-generic.  Then every \(M\)-cardinal remains a cardinal of
\(M[G]\), and every ordinal has the same cofinality as computed in \(M\) and
\(M[G]\).
\end{hbtheorem}

\begin{hbdefinition}{def:p06-nice-name}{Nice name for a subset of \(\omega\)}
The adjective \emph{nice} isolates the antichain data needed to count names;
an arbitrary name does not display that data.  A \emph{nice
\(\mathbb P\)-name for a subset of \(\omega\)} is
\[
  \dot x=\bigcup_{n<\omega}
  \{\langle\check n,p\rangle:p\in A_n\},
\]
where every \(A_n\) is an antichain in \(\mathbb P\).  This is a
metalinguistic subclass of the names in
Section~\ref{sec:constructibility-forcing-02}; it does not change evaluation.  Its
nonredundant use is that ccc makes each \(A_n\) countable.
\index{nice name!forcing}
\end{hbdefinition}

\begin{hbtheorem}{thm:p06-nice-name-counting}{proved=annex:p6.force.10}{Nice-name reduction and counting}
Under Convention~\ref{def:p06-ccc-forcing-standing}, assume \(M\) regards
\(\mathbb P\) as ccc.  If \(p\in\mathbb P\) and
\(p\forces\nu\subseteq\check\omega\), then some nice
\(\mathbb P\)-name \(\zeta\) satisfies \(p\forces\nu=\zeta\).
Consequently, if \(\mathbb P\) is infinite and
\(\lambda=(\abs{\mathbb P}^{\aleph_0})^M\), then
\[
  M[G]\entails 2^{\aleph_0}\leq\lambda.
\]
\end{hbtheorem}

\begin{hbcorollary}{thm:p06-many-cohen-continuum}{proved=annex:p6.force.11}{Forcing the size of the continuum}
Under Convention~\ref{def:p06-ccc-forcing-standing}, suppose \(M\) regards
\(\kappa\) as an infinite cardinal and satisfies
\(\kappa^{\aleph_0}=\kappa\), and let
\(\mathbb P=\operatorname{Add}(\omega,\kappa)^M\).  For every
\(M\)-generic \(G\),
\[
  M[G]\entails 2^{\aleph_0}=\kappa.
\]
In particular, a ground model of \(\mathsf{GCH}\) with
\(\kappa=\aleph_2\) yields a forcing extension satisfying
\(\neg\mathsf{CH}\).
\end{hbcorollary}

% END HISTORICAL SLICE 324-372
